% Chen-Yu Kuo — CV (English).  Build: tectonic -X compile cv-en.tex
\documentclass[10pt,a4paper]{article}

\usepackage{fontspec}
\setmainfont{texgyrepagella}[Extension=.otf,
  UprightFont=*-regular, BoldFont=*-bold, ItalicFont=*-italic, BoldItalicFont=*-bolditalic]
\setsansfont{texgyreheros}[Extension=.otf,
  UprightFont=*-regular, BoldFont=*-bold, ItalicFont=*-italic, BoldItalicFont=*-bolditalic]
\usepackage{xeCJK}
\setCJKmainfont{Heiti TC}
\XeTeXlinebreaklocale "zh"
% Keep Western punctuation in the Latin font (xeCJK otherwise renders
% em/en dashes and quotes with the CJK face as full-width glyphs).
\xeCJKDeclareCharClass{Default}{"2013, "2014, "2018, "2019, "201C, "201D, "2026, "2032}

\newcommand{\cvfootline}{Chen-Yu Kuo \textperiodcentered\ Curriculum Vitae}
\usepackage{cvstyle}

\begin{document}

\cvname{Chen-Yu Kuo}{郭宸毓}
\cvtagline{Systems and full-stack engineer — real-time earthquake early warning, from sensor firmware to the app that rings}
\cvcontact{%
  \faMapMarker*\enspace Kaohsiung, Taiwan \dt
  \faEnvelope\enspace \href{mailto:whes1015@gmail.com}{whes1015@gmail.com} \dt
  \faGithub\enspace \href{https://github.com/whes1015}{github.com/whes1015} \dt
  \faRobot\enspace \href{https://huggingface.co/YuYu1015}{huggingface.co/YuYu1015} \dt
  \faGlobe\enspace \href{https://exptech.dev}{exptech.dev}}

\vspace{7pt}
{\fontsize{9.4}{12.4}\selectfont
I founded and lead \textbf{ExpTech}, a 29-person volunteer engineering organisation that builds and operates Taiwan's
largest independent earthquake early-warning platform. I work the whole path a seismic alert travels: the ESP32
accelerometer firmware I wrote register by register, the multi-region ingest and alerting fleet it reports to, the
ground-motion models that decide what the alert says, and the Flutter app that over 100,000 people installed to receive it.
\par}

% ------------------------------------------------------------------ experience
\section{Experience}

\role{Founder \& Lead Engineer}{ExpTech (探索科技)}{Taiwan — volunteer engineering organisation, 29 members, 234 repositories}{Oct 2020 — Present}
\begin{itemize}
  \item Lead the organisation as one of its two administrators. Authored \num{13,411} commits inside it
        (\num{19,284} across GitHub, \num{477} pull requests) and am the top contributor on every flagship codebase:
        DPIP \num{838}/2,339, TREM-Lite \num{448}/847, platform API \num{480}/803 — and sole author of the
        665-commit production infrastructure repository.
  \item Signed a data agreement with Taiwan's \textbf{Central Weather Administration (CWA)} in late 2022; reported by
        \textit{親子天下} as the Seismological Center's first non-corporate partner. Invited back to the centre
        10+ times as a collaborator (\textit{天下雜誌}, Feb 2026).
  \item Ship under real stakes to real users: \textbf{100,000+} Google Play installs and \textbf{4.2/5} from 1,378 App Store ratings for DPIP; \textbf{328,542} installer downloads across TREM desktop releases.
        During the M7.2 Hualien earthquake (3 Apr 2024) the network recorded 800+ events in a single day.
  \item Grew the codebase from a school-project origin into a fleet spanning firmware, 37 Go services, Flutter
        clients and ML pipelines — recruiting and reviewing a distributed team of student engineers.
\end{itemize}

% ------------------------------------------------------------------ projects
\section{Selected Engineering Work}

\project{Ground-motion \& phase-picking models}{Python \textperiodcentered\ PyTorch \textperiodcentered\ SeisBench \textperiodcentered\ ONNX \textperiodcentered\ sole author}{huggingface.co/YuYu1015}{https://huggingface.co/YuYu1015}
\begin{itemize}
  \item Refit Taiwan's shaking-intensity model on \num{45,398} CWA PGA/PGV/intensity observations (2000–2025),
        replacing a legacy equation that over-amplified with magnitude. Added a direct PGV ground-motion prediction
        equation and corrected a site-amplification double-count that inflated estimates roughly 2$\times$.
        Held-out result: \textbf{62.0\% exact intensity vs 39.7\%} for the model then in production,
        MAE \textbf{0.400 vs 0.671}, with 0\% false alarms. Released as \textit{IntensityGMM-TW-45k-v1}.
  \item Trained a 50\,Hz PhaseNet P/S phase picker — a 6,457-line pipeline covering dataset ingest, noise injection,
        evaluation and ONNX export. Diagnosed a windowing bug in which the source window was shorter than the model
        window, teaching the model that P always sits mid-window and emitting a false pick every 45\,s stride.
        Released as \textit{PhaseNet-50Hz-268k-v1}.
  \item Built a P-window $\rightarrow$ S-amplitude regressor: 1–5\,s post-P features ($\tau_c$ after
        Wu \& Kanamori 2005, spectral centroid and bandwidth, kurtosis, zero crossings) into CNN and tabular models.
\end{itemize}

\project{ES-Net — seismometer firmware \& national sensor network}{C++ \textperiodcentered\ ESP32 \textperiodcentered\ FreeRTOS \textperiodcentered\ sole author, 7,759 LOC}{}{}
\begin{itemize}
  \item Hand-wrote an \textbf{LSM6DS3 accelerometer driver at register level} — own register map, \textsc{who\_am\_i},
        \textsc{ctrl1\_xl}, \textsc{int1\_ctrl} and FIFO paths — rather than depending on a vendor library.
  \item On-device \textbf{JMA seismic-intensity} computation: an IIR filter bank plus a binary indexed tree for the
        running amplitude percentile the JMA scale requires.
  \item Timing is the hard problem in a distributed seismometer: SNTP using \textsc{smooth} sync mode so the clock
        slews rather than steps mid-waveform, on a 60\,s interval — alongside hardware and software watchdogs,
        OTA update, and a documented custom binary wire protocol with TOTP station authentication.
  \item Feeds \textbf{200+ TREM-Net stations} nationwide; the first-generation units were self-assembled for under
        NT\$800 each.
\end{itemize}

\project{TREM — real-time earthquake monitoring desktop platform}{JavaScript/Electron \textperiodcentered\ Rust \textperiodcentered\ 93 stars \textperiodcentered\ top contributor 448/847}{ExpTechTW/TREM-Lite}{https://github.com/ExpTechTW/TREM-Lite}
\begin{itemize}
  \item Desktop client that renders live ground motion from TREM-Net and CWA feeds. \textbf{328,542 installer
        downloads} across 74 releases of the current and previous generations, shipped continuously since Jan 2023.
  \item Designed the plugin system third parties extend the client through and maintain its registry (417 commits);
        wrote the companion \texttt{trem-monitor} station-health service in Rust (105/150 commits).
\end{itemize}

\project{QZSS L1S disaster-message decoder}{C++ \textperiodcentered\ sole author \textperiodcentered\ {\raise0.06ex\hbox{$\sim$}}12k LOC with unit tests}{}{}
\begin{itemize}
  \item Decodes QZSS L1S 250-bit SBAS frames end to end — preamble, 6-bit message type, 212-bit payload and
        CRC-24Q — including type 43 \textbf{DC Report} (JMA disaster broadcasts, IS-QZSS-DCR) and type 44 DCX.
\end{itemize}

\project{Production infrastructure}{Go \textperiodcentered\ eBPF/XDP \textperiodcentered\ OpenResty \textperiodcentered\ CoreDNS \textperiodcentered\ sole author}{}{}
\begin{itemize}
  \item \textbf{kekkai} — packet filter in Go using \texttt{cilium/ebpf}, with an XDP ingress program and a TC egress
        hook; \textsc{lru\_hash}/\textsc{array}/\textsc{ringbuf} maps drive per-IP rate limiting, allow/blocklists and
        port policy, with an emergency detach path.
  \item \textbf{Nginx/OpenResty fleet} — 23 upstreams across nodes in Taipei, Kaohsiung, Tainan, Tokyo and Osaka,
        with Lua modules for traffic accounting, NTP and image proxying.
  \item \textbf{Split-horizon CoreDNS} with a custom Go \texttt{dnsstats} plugin (\texttt{miekg/dns}) exposing
        per-domain and per-client query statistics plus Prometheus metrics.
  \item Sole author too of the platform's NTP server, API gateway and Meshtastic LoRa ingest service — among the
        37 Go services the organisation runs.
\end{itemize}

\project{DPIP — cross-platform disaster alerting}{Dart/Flutter \textperiodcentered\ TypeScript \textperiodcentered\ top contributor}{ExpTechTW/DPIP}{https://github.com/ExpTechTW/DPIP}
\begin{itemize}
  \item iOS and Android client integrating TREM-Net and CWA feeds with push alerting, MapLibre layers and
        10 interface languages. Wrote the repository's test suite and its custom commit-gate tooling outright,
        and authored the backend service \texttt{dpip-server-ts} in full.
\end{itemize}

\project{LLM inference \& model compression}{PyTorch \textperiodcentered\ vLLM \textperiodcentered\ llm-compressor \textperiodcentered\ 27 models, {\raise0.06ex\hbox{$\sim$}}10,600 downloads}{}{}
\begin{itemize}
  \item Quantise large open models to \textbf{NVFP4} (TensorRT Model Optimizer), INT4-AutoRound and GGUF for
        vLLM and llama.cpp serving, including MoE architectures up to 35B.
  \item Implemented \textbf{refusal-direction ablation} — deriving the refusal direction from paired prompts and
        orthogonalising it out of the weight matrices — extending the upstream reference implementation, and
        released the \textit{Ornith-1.0} 9B/35B line including a DPO-tuned variant.
\end{itemize}

\project{Security \& incident response}{honeypot \textperiodcentered\ malware triage}{}{}
\begin{itemize}
  \item Operated a honeypot that captured \textbf{in-the-wild exploitation of CVE-2025-55182} (React Server
        Components, CVSS 10.0) within days of its December 2025 disclosure; documented six malware samples with
        verified hashes and fifteen attacker source addresses.
\end{itemize}

% ------------------------------------------------------------------ skills
\section{Technical Skills}
\skl{Languages}{TypeScript/JavaScript, Go, Python, Dart, C/C++, Rust, Lua, Shell}
\skl{Systems}{Linux, eBPF/XDP, Docker, Nginx/OpenResty, CoreDNS, WebSocket \& SSE, NTP, Prometheus, CI/CD}
\skl{Embedded}{ESP32/Arduino, FreeRTOS, I\textsuperscript{2}C/SPI, LSM6DS3, GNSS (QZSS L1S), LoRa/Meshtastic, RTL-SDR}
\skl{ML}{PyTorch, SeisBench, ONNX, scikit-learn, vLLM, llm-compressor, TensorRT Model Optimizer}
\skl{Domain}{Earthquake early warning, GMPE calibration, CWA/JMA intensity scales, phase picking, NonLinLoc travel-time grids}

% ------------------------------------------------------------------ education
\section{Education}
% TODO: add your expected graduation, e.g. {Expected Jun 2027} in the last argument.
\role{B.S., Electrical Engineering}{National Kaohsiung University of Science and Technology (NKUST)}{}{Kaohsiung, Taiwan}

% ------------------------------------------------------------------ press
\section{Open Source \& Community}
\begin{itemize}
  \item \textbf{12 merged pull requests} to projects outside my own organisation — a rendering-performance fix to
        \textbf{AnyShake}'s \texttt{spectrogram-js} (an open-source seismograph suite), a Leaflet~$\rightarrow$~MapLibre
        migration for \texttt{smc.peering.tw}, and Flutter toolchain upgrades to \textbf{NKUST-AP-Flutter},
        my university's official student app.
\end{itemize}

\section{Press \& Recognition}
\begin{itemize}
  \item \textbf{天下雜誌} (CommonWealth Magazine), Feb 2026 — quoted by name as an ExpTech collaborator in a profile
        of the CWA Seismological Center.
  \item \textbf{親子天下 翻轉教育}, Jul 2024 — feature on the DPIP team; the app reached roughly 500,000 downloads
        following the 3 April 2024 earthquake.
\end{itemize}

\note{Provenance — commit counts are from GitHub's commit-search index and per-repository contributor statistics
(retrieved Sep 2026). The TREM download figure counts installer assets only (\texttt{.exe}/\texttt{.dmg}/\texttt{.deb}/\texttt{.AppImage}/\texttt{.zip}),
excluding 4.3M auto-updater metadata requests that GitHub also reports as downloads.}

\end{document}
